% part: intuitionistic-logic
% chapter: soundness-completeness
% section: canonical-model

\documentclass[../../../include/open-logic-section]{subfiles}

\begin{document}

\olfileid{int}{sc}{mod}

\olsection{مدل کانونی}

جهان‌های مدل ما دنباله‌های متناهی~$\sigma$ از اعداد طبیعی خواهند بود؛
یعنی~$\sigma \in \Nat^*$. توجه کنید که~$\Nat^*$ به‌طور استقرایی چنین
تعریف می‌شود:
\begin{enumerate}
\item $\emptyseq \in \Nat^*$.
\item اگر~$\sigma \in \Nat^*$ و~$n \in \Nat$، آنگاه
  $\sigma.n \in \Nat^*$ (که در آن~$\sigma.n$ همان
  $\sigma \concat \tuple{n}$ است و~$\sigma \concat \sigma'$ الحاق
  $\sigma$ و~$\sigma'$ است).
\item هیچ چیز دیگری در~$\Nat^*$ نیست.
\end{enumerate}
پس می‌توانیم از~$\Nat^*$ برای ارائهٔ تعریف‌های استقرایی استفاده کنیم.

فرض کنید~$\tuple{!B_1, !C_1}$، $\tuple{!B_2, !C_2}$، \dots، شمارشی
از همهٔ جفت‌های !!{formula}ها باشد. با داشتن مجموعهٔ
!!{formula}های~$\Delta$، $\Delta(\sigma)$ را با استقرا چنین تعریف کنید:
\begin{enumerate}
\item $\Delta(\emptyseq) = \Delta$
\item $\Delta(\sigma.n) = {}$
  \[
  \begin{cases}
    (\Delta(\sigma) \cup \{!B_n\})^* &
    \text{اگر~$\Delta(\sigma) \cup \{!B_n\} \Proves/ !C_n$} \\
    \Delta(\sigma) & \text{در غیر این صورت}
  \end{cases}
  \]
\end{enumerate}
در اینجا مقصود از~$(\Delta(\sigma) \cup \{!B_n\})^*$ مجموعهٔ اولی از
!!{formula}هاست که با اعمال~\olref[lin]{lem:lindenbaum} بر مجموعهٔ
$\Delta(\sigma) \cup \{!B_n\}$ و !!{formula}~$!C_n$ وجودش تضمین
می‌شود. توجه کنید بنا بر این تعریف، اگر
$\Delta(\sigma) \cup \{!B_n\} \Proves/ !C_n$، آنگاه
$\Delta(\sigma.n) \Proves !B_n$ و~$\Delta(\sigma.n) \Proves/ !C_n$.
همچنین برای هر~$n$، $\Delta(\sigma) \subseteq \Delta(\sigma.n)$.
اگر~$\Delta$ اول باشد، آنگاه~$\Delta(\sigma)$ برای هر~$\sigma$ اول است.

\begin{defn}\ollabel{defn:canonical-model}
  فرض کنید~$\Delta$ اول باشد. \emph{مدل کانونی}~$\mModel{M(\Delta)}$
  برای~$\Delta$ چنین تعریف می‌شود:
  \begin{enumerate}
  \item $W = \Nat^*$، مجموعهٔ دنباله‌های متناهی اعداد طبیعی.
  \item $R$ ترتیب جزئی‌ای است که بنا بر آن~$R\sigma\sigma'$ اگر و تنها
    اگر~$\sigma$ قطعهٔ آغازین~$\sigma'$ باشد (یعنی برای دنباله‌ای
    $\sigma''$، $\sigma' = \sigma \concat \sigma''$).
  \item $V(p) = \Setabs{\sigma}{p \in \Delta(\sigma)}$.
  \end{enumerate}
\end{defn}

به‌آسانی می‌توان وارسی کرد که~$R$ واقعاً یک ترتیب جزئی است. شرط یکنوایی
روی~$V$ نیز برآورده است. چون~$\Delta(\sigma) \subseteq
\Delta(\sigma.n)$، با استقرا بر~$\sigma$ نتیجه می‌گیریم هرگاه
$R\sigma\sigma'$، آنگاه~$\Delta(\sigma) \subseteq \Delta(\sigma')$.

\end{document}
